\documentclass[10 pt]{article}
\usepackage[utf8]{inputenc}
\usepackage{parskip}
\usepackage[utf8]{inputenc}
\usepackage[spanish]{babel}
\usepackage[left=2cm, right=2cm, top=1.5cm, bottom=2cm]{geometry}
\usepackage{float}
\usepackage{graphicx}
\usepackage{caption}
\usepackage{cancel}
\usepackage{amsmath,amsthm,amsfonts,amscd,amssymb,amsbsy}

\title{\textbf{\underline{Los enteros}}}
\author{}
\date{}

\begin{document}

\pagestyle{empty}
\maketitle
\thispagestyle{empty}

\section*{}

\LARGE{\textbf{\underline{Escritura e identificación}}}\\ \\

\Large{Hasta ahora hemos trabajado con los números enteros mediante representantes de sus respectivas clases de equivalencia. La suma y la multiplicación se definieron operando sobre dichos representantes, y comprobamos que el resultado no depende del representante escogido. Esto sugiere la posibilidad de elegir, para cada entero, un ``representante distinguido'' o ``canónico'':} \\

\begin{itemize}
    \item \Large{\textbf{Teorema:} Sea $[(a,b)] \in \mathbb Z$. Si $[(a,b)] \neq [(0,0)]$, entonces existe $n \in \mathbb N \setminus \{0\}$ tal que $[(a,b)]=[(n,0)]$ o $[(a,b)]=[(0,n)]$.} \\
\end{itemize}

\Large{\textbf{Demostración:} Como $[(a,b)] \neq [(0,0)]$, resulta que $a+0 \neq 0+b$. De modo pues que $a \neq b$. Luego, por tricotomía, $a<b$ o $b<a$. Consideremos ambos casos:} \\

\Large{($\ast$) Supongamos que $a<b$. Entonces, como se demostró anteriormente, existe $p \in \mathbb N \setminus \{0\}$ tal que $b=a+p$. Con lo cual:}
\begin{align*}
    b&=b \\
    a+p&=b+0 \\
    &=0+b
\end{align*}
\Large{Por lo tanto, $[(a,b)]=[(0,p)]$.} \\

\Large{($\ast$) Similarmente: supongamos que $b<a$. Entonces, como se demostró anteriormente, existe $p \in \mathbb N \setminus \{0\}$ tal que $a=b+p$. Con lo cual:}
\begin{align*}
    a&=a \\
    a+0&=b+p \\
    &=p+b
\end{align*}
\Large{Por lo tanto, $[(a,b)]=[(p,0)]$.} \\
\begin{flushright}
   $\square$
\end{flushright}

\Large{El resultado anterior nos acerca a la representación habitual de los números enteros. El siguiente paso consiste en explicar cómo los números naturales se identifican naturalmente con un subconjunto de $\mathbb Z$, como cabe esperar de una construcción cuyo propósito es extender el conjunto de los números naturales.} \\

\begin{itemize}
    \item \Large{\textbf{Definición:} La siguiente función se conoce como \textbf{inyección canónica} de $\mathbb N$ en $\mathbb Z$:} 
    \begin{align*}
        i:\mathbb N &\rightarrow \mathbb Z \\
        n &\mapsto [(n,0)]
    \end{align*}
\end{itemize}
\begin{flushright}
    $\blacksquare$
\end{flushright}

\Large{El papel de la inyección canónica es identificar de manera natural a $\mathbb N$ con un subconjunto de $\mathbb Z$. Para comprobar que esta identificación es adecuada, verificaremos que la aplicación es inyectiva y que preserva la suma y la multiplicación:} \\

\begin{itemize}
    \item \Large{\textbf{Proposición:} La inyección canónica de $\mathbb N$ en $\mathbb Z$ es inyectiva.} \\
\end{itemize}

\Large{\textbf{Demostración:} Sean $n,m \in \mathbb N$, y supongamos que $i(n)=i(m)$. Es decir que $[(n,0)]=[(m,0)]$, y por tanto que $n+0=0+m$. Con lo cual, $n=m$. Tenemos entonces que $i$ es inyectiva.} \\

\begin{flushright}
    $\square$
\end{flushright}

\Large{La proposición anterior garantiza que podemos identificar los números naturales con su imagen en $\mathbb Z$ sin perder información: \textit{cada número natural queda representado por un único número entero}. Sin embargo, para que esta identificación respete la estructura algebraica de ambos conjuntos, resta comprobar que la inyección canónica preserva la suma y la multiplicación:} \\

\begin{itemize}
    \item \Large{\textbf{Proposición:} Sean $n,m \in \mathbb N$. \\

    i) $i(n+m)=i(n)+i(m)$. \\

    ii) $i(n \cdot m)=i(n) \cdot i(m)$.} \\
\end{itemize}

\Large{\textbf{Demostración:} \\

i) Observe que:}
\begin{align*}
    i(n+m)&=[(n+m,0)] \\
    &=[(n,0)]+[(m,0)] \\
    &=i(n)+i(m) 
\end{align*}
\Large{Observe que:}
\begin{align*}
    i(n \cdot m)&=[(n \cdot m,0)] \\
    &=[(n \cdot m,0 \cdot 0,n \cdot0+0 \cdot m)] \\
    &=[(n,0)] \cdot [(m,0)] \\
    &=i(n) \cdot i(m)
\end{align*}

\begin{flushright}
    $\square$
\end{flushright}

\Large{Como se anticipó, la inyección canónica nos permite definir un subconjunto de $\mathbb Z$ que reproduce la estructura algebraica de $\mathbb N$:} \\

\begin{itemize}
    \item \Large{\textbf{Definición:} $\mathbb N_1:=i(\mathbb N)$. Es decir que: \\

    -- $\mathbb N_1=\{x \in \mathbb Z \hspace{0.2 cm}|\hspace{0.2 cm}(\exists n \in \mathbb N)(x=[(n,0)])\}$.} \\
\end{itemize}

\begin{flushright}
    $\blacksquare$
\end{flushright}

\Large{La preservación de las operaciones permite trasladar inmediatamente a $\mathbb N_1$ cualquier propiedad algebraica de $\mathbb N$. Veamos un ejemplo:} \\

\begin{itemize}
    \item \Large{\textbf{Proposición:} Sean $n,m \in \mathbb N$. Si $i(n)+i(m)=[(0,0)]$, entonces $i(n)=i(m)=[(0,0)]$.} \\
\end{itemize}

\Large{\textbf{Demostración:} Note que $i(0)=[(0,0)]$. Con lo cual:}
\begin{align*}
    i(n)+i(m)&=[(0,0)] \\
    i(n+m)&=i(0)
\end{align*}
\Large{Como $i$ es inyectiva, se sigue que $n+m=0$. Luego, como se demostró anteriormente, $n=m=0$. Por lo tanto, $i(n)=i(m)=i(0)=[(0,0)]$.} \\

\begin{flushright}
    $\square$
\end{flushright}

\Large{Este fenómeno no es exclusivo de la propiedad anterior. En general, cualquier propiedad algebraica de los números naturales expresada en términos de la suma, la multiplicación y la igualdad puede trasladarse a $\mathbb N_1$ mediante la inyección canónica. Del mismo modo, las demostraciones por inducción pueden realizarse indistintamente en $\mathbb N$ o en $\mathbb N_1$: demostrar un resultado en uno de estos conjuntos equivale a demostrar el correspondiente en el otro:} \\

\begin{itemize}
    \item \Large{\textbf{Teorema:} Sea $A \subseteq \mathbb N$. \\

    i) Si $A$ es inductivo, entonces $i(A)=\mathbb N_1$. \\

    ii) Si $[(0,0)] \in i(A)$ y $[(n,0)]+[(1,0)] \in i(A)$, para todo $[(n,0)] \in i(A)$, entonces $A=\mathbb N$ y $i(A)=\mathbb N_1$.} \\
\end{itemize}

\Large{\textbf{Demostración:} \\

i) De acuerdo con la hipótesis, $A=\mathbb N$ (por el principio de inducción). Luego $i(A)=i(\mathbb N)=\mathbb N_1$. \\

ii) Por hipótesis, $i(0) \in i(A)$ (porque $i(0)=[(0,0)]$). Como $i$ es inyectiva, necesariamente $0 \in A$. A continuación, consideremos $n \in A$. Como $[(n,0)] \in i(A)$ (porque $i(n) \in i(A)$), resulta que $[(n,0)]+[(1,0)] \in i(A)$. Es decir que $[(n+1,0)] \in i(A)$, y por tanto que $i(n+1) \in i(A)$. Y como $i$ es inyectiva, concluimos que $n+1 \in A$. Tenemos entonces que $A$ es inductivo, y por tanto que $A=\mathbb N$. Luego, por la parte anterior, $i(A)=\mathbb N_1$.} \\

\begin{flushright}
    $\square$
\end{flushright}

\Large{Desde el punto de vista algebraico y lógico, no existe diferencia entre trabajar en $\mathbb N$ y hacerlo en $\mathbb N_1$: trabajaremos indistintamente en uno u otro, según convenga. En consecuencia, a partir de este momento identificaremos ambos conjuntos y consideraremos a $\mathbb N$ como un subconjunto de $\mathbb Z$. Por ello, escribiremos ``$\mathbb N$'' en lugar de ``$\mathbb N_1$''. Esta identificación, junto con la existencia de representantes canónicos, nos permite describir los números enteros de una manera mucho más sencilla que la empleada durante su construcción. En efecto, cada entero puede representarse mediante un único número natural o el inverso aditivo de uno. Así pues, adoptaremos tales representantes como la forma habitual de representar a los enteros, recuperando así la notación usual:} \\

\begin{itemize}
    \item \Large{\textbf{Definición:} Sea $n \in \mathbb N$. Entonces: \\

    i) $n:=[(n,0)]$. \\

    ii) $0:=[(0,0)]$. \\

    iii) $-n:=[(0,n)]$.} \\
\end{itemize}

\Large{\textbf{Nota:} El teorema demostrado al inicio de esta sección garantiza que todo número entero puede representarse de esta forma.} \\

\begin{flushright}
    $\blacksquare$
\end{flushright}

\Large{La notación $-n$ para representar al entero $[(0,n)]$ no es arbitraria: ya demostramos que $[(0,n)]$ es el inverso aditivo del entero $n=[(n,0)]$. Por esta razón, resulta natural adoptar la notación habitual $-n=[(0,n)]$.} \\

\Large{La \textbf{resta} se define como es usual en la teoría de anillos:} \\

\begin{itemize}
    \item \Large{\textbf{Definición:} Sean $n,m \in \mathbb Z$. Entonces $n-m:=n+(-m)$. Nos referimos a $n-m$ como la \textbf{resta} de $n$ y $m$.} \\
\end{itemize}

\Large{\textbf{Observación:} Se sigue de inmediato de la definición que $n-m \in \mathbb Z$.} \\

\begin{flushright}
    $\blacksquare$
\end{flushright}

\Large{Con esta notación, cada entero puede interpretarse como una diferencia de números naturales:} \\

\begin{itemize}
    \item \Large{\textbf{Proposición:} Sean $[(a,b)] \in \mathbb Z$. Entonces $[(a,b)]=a-b$.} \\
\end{itemize}

\Large{\textbf{Demostración:} Por definición:}
\begin{align*}
    [(a,b)]&=[(a,0)]+[(0,b)] \\
    &=a+(-b) \\
    &=a-b
\end{align*}

\begin{flushright}
    $\square$
\end{flushright}

\Large{Gracias a la identificación de $\mathbb N$ con un subconjunto de $\mathbb Z$, toda pareja de números naturales puede restarse sin restricciones. El resultado pertenece siempre a $\mathbb Z$, aunque puede no pertenecer a $\mathbb N$.} \\

\end{document}